\documentclass[journal,12pt,twocolumn]{IEEEtran}
%
\usepackage{setspace}
\usepackage{gensymb}
%\doublespacing
\singlespacing

%\usepackage{graphicx}
%\usepackage{amssymb}
%\usepackage{relsize}
\usepackage[cmex10]{amsmath}
%\usepackage{amsthm}
%\interdisplaylinepenalty=2500
%\savesymbol{iint}
%\usepackage{txfonts}
%\restoresymbol{TXF}{iint}
%\usepackage{wasysym}
\usepackage{amsthm}
\usepackage{iithtlc}
\usepackage{mathrsfs}
\usepackage{txfonts}
\usepackage{stfloats}
\usepackage{bm}
\usepackage{cite}
\usepackage{cases}
\usepackage{subfig}
%\usepackage{xtab}
\usepackage{longtable}
\usepackage{multirow}
%\usepackage{algorithm}
%\usepackage{algpseudocode}
\usepackage{enumitem}
\usepackage{mathtools}
\usepackage{tikz}
\usepackage{circuitikz}
\usepackage{verbatim}
\usepackage{tfrupee}
\usepackage[breaklinks=true]{hyperref}
%\usepackage{stmaryrd}
\usepackage{tkz-euclide} % loads  TikZ and tkz-base
\usetkzobj{all}
\usepackage{listings}
    \usepackage{color}                                            %%
    \usepackage{array}                                            %%
    \usepackage{longtable}                                        %%
    \usepackage{calc}                                             %%
    \usepackage{multirow}                                         %%
    \usepackage{hhline}                                           %%
    \usepackage{ifthen}                                           %%
  %optionally (for landscape tables embedded in another document): %%
    \usepackage{lscape}     
\usepackage{multicol}
\usepackage{chngcntr}
%\usepackage{enumerate}

%\usepackage{wasysym}
%\newcounter{MYtempeqncnt}
\DeclareMathOperator*{\Res}{Res}
%\renewcommand{\baselinestretch}{2}
\renewcommand\thesection{\arabic{section}}
\renewcommand\thesubsection{\thesection.\arabic{subsection}}
\renewcommand\thesubsubsection{\thesubsection.\arabic{subsubsection}}

\renewcommand\thesectiondis{\arabic{section}}
\renewcommand\thesubsectiondis{\thesectiondis.\arabic{subsection}}
\renewcommand\thesubsubsectiondis{\thesubsectiondis.\arabic{subsubsection}}

% correct bad hyphenation here
\hyphenation{op-tical net-works semi-conduc-tor}
\def\inputGnumericTable{}                                 %%

\lstset{
%language=C,
frame=single, 
breaklines=true,
columns=fullflexible
}
%\lstset{
%language=tex,
%frame=single, 
%breaklines=true
%}

\begin{document}
%


\newtheorem{theorem}{Theorem}[section]
\newtheorem{problem}{Problem}
\newtheorem{proposition}{Proposition}[section]
\newtheorem{lemma}{Lemma}[section]
\newtheorem{corollary}[theorem]{Corollary}
\newtheorem{example}{Example}[section]
\newtheorem{definition}[problem]{Definition}
%\newtheorem{thm}{Theorem}[section] 
%\newtheorem{defn}[thm]{Definition}
%\newtheorem{algorithm}{Algorithm}[section]
%\newtheorem{cor}{Corollary}
\newcommand{\BEQA}{\begin{eqnarray}}
\newcommand{\EEQA}{\end{eqnarray}}
\newcommand{\define}{\stackrel{\triangle}{=}}

\bibliographystyle{IEEEtran}
%\bibliographystyle{ieeetr}


\providecommand{\mbf}{\mathbf}
\providecommand{\pr}[1]{\ensuremath{\Pr\left(#1\right)}}
\providecommand{\qfunc}[1]{\ensuremath{Q\left(#1\right)}}
\providecommand{\sbrak}[1]{\ensuremath{{}\left[#1\right]}}
\providecommand{\lsbrak}[1]{\ensuremath{{}\left[#1\right.}}
\providecommand{\rsbrak}[1]{\ensuremath{{}\left.#1\right]}}
\providecommand{\brak}[1]{\ensuremath{\left(#1\right)}}
\providecommand{\lbrak}[1]{\ensuremath{\left(#1\right.}}
\providecommand{\rbrak}[1]{\ensuremath{\left.#1\right)}}
\providecommand{\cbrak}[1]{\ensuremath{\left\{#1\right\}}}
\providecommand{\lcbrak}[1]{\ensuremath{\left\{#1\right.}}
\providecommand{\rcbrak}[1]{\ensuremath{\left.#1\right\}}}
\theoremstyle{remark}
\newtheorem{rem}{Remark}
\newcommand{\sgn}{\mathop{\mathrm{sgn}}}
\providecommand{\abs}[1]{\left\vert#1\right\vert}
\providecommand{\res}[1]{\Res\displaylimits_{#1}} 
\providecommand{\norm}[1]{\left\lVert#1\right\rVert}
%\providecommand{\norm}[1]{\lVert#1\rVert}
\providecommand{\mtx}[1]{\mathbf{#1}}
\providecommand{\mean}[1]{E\left[ #1 \right]}
\providecommand{\fourier}{\overset{\mathcal{F}}{ \rightleftharpoons}}
%\providecommand{\hilbert}{\overset{\mathcal{H}}{ \rightleftharpoons}}
\providecommand{\system}{\overset{\mathcal{H}}{ \longleftrightarrow}}
	%\newcommand{\solution}[2]{\textbf{Solution:}{#1}}
\newcommand{\solution}{\noindent \textbf{Solution: }}
\newcommand{\cosec}{\,\text{cosec}\,}
\providecommand{\dec}[2]{\ensuremath{\overset{#1}{\underset{#2}{\gtrless}}}}
\newcommand{\myvec}[1]{\ensuremath{\begin{pmatrix}#1\end{pmatrix}}}
\newcommand{\mydet}[1]{\ensuremath{\begin{vmatrix}#1\end{vmatrix}}}
%\numberwithin{equation}{section}
\numberwithin{equation}{subsection}
%\numberwithin{problem}{section}
%\numberwithin{definition}{section}
\makeatletter
\@addtoreset{figure}{problem}
\makeatother

\let\StandardTheFigure\thefigure
\let\vec\mathbf
%\renewcommand{\thefigure}{\theproblem.\arabic{figure}}
\renewcommand{\thefigure}{\theproblem}
%\setlist[enumerate,1]{before=\renewcommand\theequation{\theenumi.\arabic{equation}}
%\counterwithin{equation}{enumi}


%\renewcommand{\theequation}{\arabic{subsection}.\arabic{equation}}

\def\putbox#1#2#3{\makebox[0in][l]{\makebox[#1][l]{}\raisebox{\baselineskip}[0in][0in]{\raisebox{#2}[0in][0in]{#3}}}}
     \def\rightbox#1{\makebox[0in][r]{#1}}
     \def\centbox#1{\makebox[0in]{#1}}
     \def\topbox#1{\raisebox{-\baselineskip}[0in][0in]{#1}}
     \def\midbox#1{\raisebox{-0.5\baselineskip}[0in][0in]{#1}}

\vspace{3cm}

\title{
	\logo{
Algebra: Maths Olympiad
	}
}
\author{ G V V Sharma$^{*}$% <-this % stops a space
	\thanks{*The author is with the Department
		of Electrical Engineering, Indian Institute of Technology, Hyderabad
		502285 India e-mail:  gadepall@iith.ac.in. All content in this manual is released under GNU GPL.  Free and open source.}
	
}	
%\title{
%	\logo{Matrix Analysis through Octave}{\begin{center}\includegraphics[scale=.24]{tlc}\end{center}}{}{HAMDSP}
%}


% paper title
% can use linebreaks \\ within to get better formatting as desired
%\title{Matrix Analysis through Octave}
%
%
% author names and IEEE memberships
% note positions of commas and nonbreaking spaces ( ~ ) LaTeX will not break
% a structure at a ~ so this keeps an author's name from being broken across
% two lines.
% use \thanks{} to gain access to the first footnote area
% a separate \thanks must be used for each paragraph as LaTeX2e's \thanks
% was not built to handle multiple paragraphs
%

%\author{<-this % stops a space
%\thanks{}}
%}
% note the % following the last \IEEEmembership and also \thanks - 
% these prevent an unwanted space from occurring between the last author name
% and the end of the author line. i.e., if you had this:
% 
% \author{....lastname \thanks{...} \thanks{...} }
%                     ^------------^------------^----Do not want these spaces!
%
% a space would be appended to the last name and could cause every name on that
% line to be shifted left slightly. This is one of those "LaTeX things". For
% instance, "\textbf{A} \textbf{B}" will typeset as "A B" not "AB". To get
% "AB" then you have to do: "\textbf{A}\textbf{B}"
% \thanks is no different in this regard, so shield the last } of each \thanks
% that ends a line with a % and do not let a space in before the next \thanks.
% Spaces after \IEEEmembership other than the last one are OK (and needed) as
% you are supposed to have spaces between the names. For what it is worth,
% this is a minor point as most people would not even notice if the said evil
% space somehow managed to creep in.



% The paper headers
%\markboth{Journal of \LaTeX\ Class Files,~Vol.~6, No.~1, January~2007}%
%{Shell \MakeLowercase{\textit{et al.}}: Bare Demo of IEEEtran.cls for Journals}
% The only time the second header will appear is for the odd numbered pages
% after the title page when using the twoside option.
% 
% *** Note that you probably will NOT want to include the author's ***
% *** name in the headers of peer review papers.                   ***
% You can use \ifCLASSOPTIONpeerreview for conditional compilation here if
% you desire.




% If you want to put a publisher's ID mark on the page you can do it like
% this:
%\IEEEpubid{0000--0000/00\$00.00~\copyright~2007 IEEE}
% Remember, if you use this you must call \IEEEpubidadjcol in the second
% column for its text to clear the IEEEpubid mark.



% make the title area
\maketitle

\newpage

%\tableofcontents

\bigskip

\renewcommand{\thefigure}{\theenumi}
\renewcommand{\thetable}{\theenumi}
%\renewcommand{\theequation}{\theenumi}

%\begin{abstract}
%%\boldmath
%In this letter, an algorithm for evaluating the exact analytical bit error rate  (BER)  for the piecewise linear (PL) combiner for  multiple relays is presented. Previous results were available only for upto three relays. The algorithm is unique in the sense that  the actual mathematical expressions, that are prohibitively large, need not be explicitly obtained. The diversity gain due to multiple relays is shown through plots of the analytical BER, well supported by simulations. 
%
%\end{abstract}
% IEEEtran.cls defaults to using nonbold math in the Abstract.
% This preserves the distinction between vectors and scalars. However,
% if the journal you are submitting to favors bold math in the abstract,
% then you can use LaTeX's standard command \boldmath at the very start
% of the abstract to achieve this. Many IEEE journals frown on math
% in the abstract anyway.

% Note that keywords are not normally used for peerreview papers.
%\begin{IEEEkeywords}
%Cooperative diversity, decode and forward, piecewise linear
%\end{IEEEkeywords}



% For peer review papers, you can put extra information on the cover
% page as needed:
% \ifCLASSOPTIONpeerreview
% \begin{center} \bfseries EDICS Category: 3-BBND \end{center}
% \fi
%
% For peerreview papers, this IEEEtran command inserts a page break and
% creates the second title. It will be ignored for other modes.
%\IEEEpeerreviewmaketitle

\begin{abstract}
This book provides a collection of the international maths olympiad problems in algebra.
\end{abstract}
%Download python codes using 
%\begin{lstlisting}
%svn co https://github.com/gadepall/school/trunk/ncert/computation/codes
%\end{lstlisting}

\renewcommand{\theequation}{\theenumi}
\begin{enumerate}[label=\arabic*.,ref=\theenumi]
%\begin{enumerate}[label=\arabic*.,ref=\thesubsection.\theenumi]
\numberwithin{equation}{enumi}

\item 100 cards are numbered 1 to 100 (each card different) and placed in 3 boxes (at least one card in each box). How many ways can this be done so that if two boxes are selected and a card is taken from each, then the knowledge of their sum alone is always sufficient to identify the third box?\\

\item Can we find N divisible by just 2000 different primes, so that N divides $2^{N} + 1$? [N may be divisible by a prime power.]





\item Prove that\\
\\ $\frac{a}{\sqrt{(a^2 + 8bc)}}$+$\frac{b}{\sqrt{(b^2 + 8ac)}}$+$\frac{c}{\sqrt{(c^2 + 8ba)}}$ $\geq 1$\\
\\ for all positive real numbers a,b and c.\\

\item Let a, b, c, d be integers with $a > b > c > d > 0$. Suppose that\\
\\ ac + bd = (b + d  a - c) (b + d - a + c).\\
\\ Prove that ab + cd is not prime.










\item S is the set of all (h, k) with h, k non-negative integers such that h + $k < n$. Each element of S is colored red or blue, so that if (h, k) is red and $h^{'}\leq h, k^{'} \leq k$, then ($h^{'}, k^{'}$) is also red. A type 1 subset of S has n blue elements with different first member and a type 2 subset of S has n blue elements with different second member. Show that there are the same number of type 1 and type 2 subsets.\\

\item Find all pairs of integers $m > 2$, $n > 2$ such that there are infinitely many positive integers k for which $k^n + k^2 - 1$ divides $k^m + k + 1$.\\

\item The positive divisors of the integer $n > 1$ are $d_1 < d_2 < ... < d_k$, so that $d_1 = 1$, $d_k = n$. Let $d = d_1d_2 + d_2d_3 + ... + d_{k-1}d_k$. Show that $d < n^2$ and find all n for which d divides $n^2$.



\item A convex hexagon has the property that for any pair of opposite sides the distance between their midpoints is $\sqrt{3}/2$ times the sum of their lengths Show that all the hexagon’s angles are equal.\\

\item ABCD is cyclic. The feet of the perpendicular from D to the lines AB, BC, CA are P, Q, R respectively. Show that the angle bisectors of ABC and CDA meet on the line AC if RP=RQ.


\item Consider in plane circle $\Gamma$ with center O and a line l not intersecting circle $\Gamma$. Prove that there is a unique point Q on the perpendicular drawn from O to the line l, such that for any point P on the line l, PQ represents the length of the tangent from P to the circle $\Gamma$.

\item Let ABCD be a quadrilateral X and Y be the mid points of AC and BD respectively and the lines through X and Y respectively parllel to BD, AC meet in O. Let P, Q, R, S be the mid points of AB, BC, CD, DA respectively. Prove that
\begin{enumerate}
\item quadrilaterals APOS and APXS have the same area.
\item the areeas is the quadrilateral APOS, BQOP, CROQ, DSOR are all equal. 
\end{enumerate}

\item Let x, y, z be three positive reals such that xyz $ \geq 1$. Prove that
\begin{align*}
\frac{x^5-x^2}{x^5 + y^2 + z^2}+\frac{y^5-y^2}{x^2 + y^5 + z^2}+\frac{z^5-z^2}{x^2 + y^2 + z^5}\geq 0
\end{align*}

\item Let $a_1, a_2,$....... be a sequence of integers with infinitely many positive and negative terms. Suppose that for every positive integer n the
numbers $a_1, a_2,........, a_n$ leave n different remainders upon division by n.
Prove that every integer occurs exactly once in the sequence $a_1, a_2,$.......

\item Determine all positive integers relatively prime to all the terms of the infinite sequence
\begin{align*}
a_n = 2^n + 3^n + 6^{n}-1, n\geq 1.
\end{align*}

\item Find the least possible value of a + b, where a, b are positive integers such that 11 divides a + 13b and 13 divides a+11b.

\item A 6 $\times$ 6 square is dissected in to 9 rectangles by lines parallel to its sides such that all these rectangles have integer sides, Prove that there are always two congruent rectangles.

\item Let ABCD be a quadrilateral in which AB is parallel to CD and perpendicular to AD ; AB = 3CD;and the area of the quadrilateral is 4. If a circle can be drawn touching all the sides of the quadrilateral, find its radius.

\item Let a, b, c be the three natural numbers such that $a < b < c$ and ged(c-a, c-b) = 1. Suppose there exists an integer d such that a+d, b+d, c+d from the sides of a right angled triangle. Prove that there exist integers l, m such that 
c + d = $l^2 + m^2$.

\item Find all pairs (a, b) of real numbers such that whenever $\alpha$ is a root of
\begin{align} 
x^2 + ax + b = 0,
\end{align}
$\alpha^2$ - 2 is also a root of the equation.


\item All the points in the plane are coloured using three colours, Prove that there exists a triangle with vertices having the same colour such that either it is isosceles or its angles are in geometric progression.
 







\item Let ABCD be a quadrilateral inscribed in a circle. Let E, F, G, H be the midpoints of the arcs AB, BC, CD, DA of the circle. Suppose AC.BD = EG.FH. Prove that AC, BD, EG, FH are concurrent.

\item Let D,E,F be points on the sides BC, CA, AB respectively of a triangle ABC such that BD = CE = AF and
 $\angle$BDF = $\angle$CED = $\angle$AFE. Prove that ABC is equilateral.






















\item Let ABC be a triangle and D be a point on the segment BC such that DC = 2BD. Let E be the mid-point of AC. Let AD and BE intersect in P. Determine the ratios BP/PE and AP/PD.

\item Let ABC be a triangle. Let BE and CF be internal angle bisector of $\angle$B and $\angle$C respectively with E on AC and F on AB. Suppose X is a point on the segment CF such that AX $\perp$ CF;and Y is a point on the segment BE such that AY $\perp$ BE. Prove that XY = (b + c - a)/2 where BC = a, CA = b and AB = c.


\item Let AD and BC be the parallel sides of a trapezium ABCD. Let P and Q be the midpoints of the diagonals AC and BD. If $AD = 16$ and $BC = 20$, What is the length of PQ?

\item In a triangle ABC, let H, I and O be the orthocentre, incentre and circumcentre, respectively. If the points B, H, I, C lie on a circle, What is magnitude of $\angle BOC$ in degrees?

\item Let ABC be an equilateral triangle. Let P and S be points on AB and AC,respectively, and let Q and R be points on BC such that PQRS is a triangle. If $PQ = \sqrt{3}PS$ and the area of PQRS is $28\sqrt{3}$, What is the length of PC?

\item Let $A_1, B_1, C_1, D_1$ be the midpoints of the sides of a convex quadrilateral ABCD and let $A_1, B_2, C_2, D_2$ be the midpoints of the quadrilateral $A_1B_1C_1D_1$. If $A_2B_2C_2D_2$ is a rectangle with sides 4 and 6, then what is the product of the lengths of the diagonals of ABCD?

\item Let S be a circle with centre O. A chord AB, not a diameter, divides S into two regions $R_1$ and $R_2$ such that O belongs to $R_2$. Let $S_1$ be a circle with centre in $R_1$, touching AB at X and S internally. Let $S_2$ be a circle with centre in $R_2$, touching AB at Y, the circle S internally and passing through the centre of S. The point X lies on the diameter passing through the centre of  $S_2$ and $\angle YXO = 30^{o}$. If the radius of $S_2$ is 100 then what is the radius of $S_1$?

\item In a triangle ABC with $\angle BCA = 90^{o}$, the perpendicular bisector of AB intersects segments AB and AC at X and Y, respectively. If the ratio of the area of quadrilateral BXYC to the area of triangle ABC is 13 : 18 and BC = 12 then what is the length of AC?

\item The first term of a sequence is 2014. Each succeeding term is the sum of the cubes of the digits of the previous term. What is the $2014^{th}$ term of the sequence?

\item In a triangle with integer side lengths, one side is three times as long as second side, and the length of the third side is 17. What is the greatest possible perimeter of the triangle?

\item For how many natural numbers n between 1 and 2014(both inclusive) is $\frac{8n}{9999 - n}$ an integer?

\item One morning, each member of Manjul's family drank an 8-ounce mixture of coffee and milk. The amounts of coffee and milk varied from cup to cup, but were never zero. Manjul drank 1/7-th of the total amount of milk and 2/17-th of the total amount of coffee. How many people are there in Manjul's family?

\item For a natural number b, let N(b) denote the number of natural numbers a for which the equation $x^2 + ax + b = 0$ has integer roots. What is the smallest value of b for which N(b)=20?

\item Let S be a set of real numbers with mean M. If the means of the sets $S \cup $ \{15\} and $S \cup $ \{15,1\} are 
M + 2 and M + 1 respectively, then how many elements does S have?

\item Let f be a one-to-one funtion from the set of natural numbers to itself such that f(mn) = f(m)f(n) for all natural numbers m and n. What is the least possible value of f(999)?

\item What is the number of ordered pairs (A, B) where A and B are subsets of \{1,2,....,5\} such that neither 
$A \subseteq B$ nor $B \subseteq A$?

\item A man walks a certain distance and rides back in $3\frac{3}{4}$ hours; he could ride both ways in $2\frac{1}{2}$ hours. How many hours would it take him to walk both ways?

\item How many line segments have both their endpoints located at the vertices of a given cube?

\item Let E(n) denote the sum of the even digits of n. For example, E(1243) = 2 + 4 = 6. What is the value of E(1) + E(2) + E(3) +.... + E(100)?

\item How many two-digit positive integers N have the property that the sum of N and the number obtained by reversing the order of the digits of N is a perfect square?

\item A 2 $\times$ 3 rectangle and a 3 $\times$ 4 rectangle are contained within a square without overlapping at any interior point, and the sides of the square are parallel to the sides of the two given rectangles. What is the smallest possible area of the square?

\item What is the greatest possible perimeter of a right-angled triangle with integer side lengths if one of the sides has length 12?

\item At a party, each man danced with exactly four women and each woman danced with exactly three men. Nine men attended the party. How many women attended the party?

\item A subset B of the set of first 100 positive integers has the property that no two elements of B sum to 125. What is the maximum possible number of elements in B?

\item The digits of a positive integer n are four consecutive integers in decreasing order when read from left to right. What is the sum of the possible remainders when n is divided by 37?

\item Let N denote the set of all natural numbers. Define a function T : N $\to$ N by
T(2k) = k and T(2k + 1) = 2k + 2. We write $T^{2}(n)$ = T(T(n)) and in general
$T^{k}(n) = T^{k-1}(T(n))$ for any $k > 1$.
\begin{enumerate}
\item Show that for each n $\in$ N, there exists k such that $T^{k}$(n) = 1.
\item For k $\in$ N, let $c_k$ denote the number of elements in the set $\{n : T^{k}(n) = 1\}$.
Prove that $c_{k+2} = c_{k+1} + c_{k}$s, for $k \geq 1$.
\end{enumerate}

\item Suppose 2016 points of the circumference of a circle are coloured red and the remaining points are coloured blue. Given any natural number n $\geq$ 3, Prove that there is a regular n-sided polygon all of whose vertices are blue.
\item Consider $n^{2}$ units squares in the xy-plane centred at point (i, j) with integer co-ordinates, 
1 $\leq$ i $\leq$ n, 1 $\leq$ j $\leq$ n. It is required to colour each unit square in such a way that when ever 
1 $\leq$ i $<$ j $\leq$ n and 1 $\leq$ k $<$ l $\leq$ n,the three squares with centres at (i, k),(j, k),(j, l)have distinct colours. What is the least possible number of colours needed?

\item Let x, y, z be real numbers each greater than 1. Prove that 
\begin{align*}
\frac{x+1}{y+1} + \frac{y+1}{z+1} + \frac{z+1}{x+1} \leq \frac{y-1}{z-1} + \frac{z-1}{x-1}.
\end{align*}

\end{enumerate}

\end{document}


